\section{Objetivos y Fundamento Teórico}

\subsection{Contexto y Objetivos del Proyecto}
El proyecto \textbf{Sistema Integral de Gestión de Capacitaciones (SIGC-CUSCO)} es una plataforma web desarrollada para registrar, gestionar y certificar cursos y talleres en entidades públicas y privadas del Cusco (Gobierno Regional, Municipalidades y DIRCETUR), eliminando el uso de planillas físicas y formatos dispersos. 

\begin{itemize}[nosep]
    \item \textbf{Objetivo General:} Diseñar la arquitectura funcional y estructural del sistema mediante diagramas UML y prototipos de interfaz de alta fidelidad, validándolos con el sistema real implementado en Laravel 13 e Inertia Vue 3.
    \item \textbf{Objetivos Específicos:}
    \begin{itemize}[nosep]
        \item Modelar el comportamiento y casos de uso con actores diferenciados y relaciones formales.
        \item Construir el diagrama de clases del dominio con atributos, métodos y multiplicidades.
        \item Diseñar e integrar prototipos para los módulos de \textbf{Login}, \textbf{Alertas} y \textbf{Reportes}.
        \item Contrastar los diseños con las pantallas del sistema en ejecución local (\texttt{localhost:8000}).
    \end{itemize}
\end{itemize}

\subsection{Herramientas Empleadas}
\begin{itemize}[nosep]
    \item \textbf{Draw.io (diagrams.net):} Modelado de diagramas en formato nativo editable (\texttt{.drawio}).
    \item \textbf{HTML5 + Tailwind CSS:} Prototipos interactivos de alta fidelidad para evaluación UX.
    \item \textbf{SIGC-CUSCO (PHP 8.3, Laravel 13, Vue 3, Vite, MySQL):} Aplicación en ejecución para extracción de evidencias y capturas de pantalla reales.
\end{itemize}
